% Faithful official CCC2022 article16 additional rKt branches used for the PH main theorem.
\section*{Additional author rKt and pseudodeterministic-PRG context}
We denote by $\mathrm{rK}_\delta^t(x)$ the minimum length of a randomized machine $M$ that outputs $x$ with probability at least $\delta$ when running for at most $t$ steps (see [38,33]). We simply write $\mathrm{rK}^t(x)$ when $\delta=2/3$.
\paragraph{Definition 8 (Pseudodeterministic PRG [39,33]).} Fix an error function $\varepsilon:\mathbb N\to\mathbb R$. A function family $G_n:\{0,1\}^\ell\to\{0,1\}^n$ is an $\varepsilon$-error Pseudorandom Generator (PRG) if, for all sufficiently large $n\in\mathbb N$, the distribution $G(U_{\ell(n)})$ $\varepsilon(n)$-fools circuits of size $n$, using the seed length $\ell=\ell(n)$, i.e., for all circuits $C$ of size $n$ on $n$ inputs,
\[\left|\Pr_{r\sim\{0,1\}^n}[C(r)=1]-\Pr_{\alpha\sim\{0,1\}^{\ell(n)}}[C(G(\alpha))=1]\right|\leq\varepsilon.\]
Such a PRG $G$ is called a pseudodeterministic PRG if there is a randomized algorithm $M_G$ such that, for every $n$ and $\alpha\in\{0,1\}^{\ell(n)}$, $M_G(1^n,\alpha)$ outputs the string $G_n(\alpha)\in\{0,1\}^n$ with probability at least $1-2^{-n}$ over its internal randomness, and $M_G$ runs in time $2^{O(\ell(n))}$.
\paragraph{Lemma 9 ([19,39]).} If BPE contains a language $L$ of circuit complexity $2^{\Omega(n)}$ for almost all input lengths $n$, then there is an $\varepsilon$-error pseudodeterministic PRG of seed length $O(\log(n/\varepsilon))$.
\paragraph{Lemma 10 ([35, Lemma 2]).} $\mathrm E^{\Sigma_2^{\mathrm P}}$ contains a language of maximum circuit complexity (at least $2^n/n$) for almost all input lengths $n$.
Using padding as in [2], we get the following.
\paragraph{Lemma 11.} If $\mathrm{DistP}^{\Sigma_2^{\mathrm P}}\subseteq\mathrm{AvgBPP}$, then $\mathrm E^{\Sigma_2^{\mathrm P}}=\mathrm{BPE}$.
\paragraph{Proof.} The inclusion $\mathrm{BPE}\subseteq\mathrm E^{\Sigma_2^{\mathrm P}}$ follows from the inclusion $\mathrm{BPP}\subseteq\Sigma_2^{\mathrm P}$ [22] by padding. For the other direction, let $L\in\mathrm E^{\Sigma_2^{\mathrm P}}$ be arbitrary. Define its padded version $L_{\mathrm{pad}}=\{(x,1^{2^{2|x|}})\mid x\in L\}$. Note that $L_{\mathrm{pad}}\in\mathrm P^{\Sigma_2^{\mathrm P}}$. Define a family of distributions $D_n=(U_n,1^{2^{2n}})$, for the uniform distribution over $n$-bit strings $U_n$. By assumption, $(L_{\mathrm{pad}},D_n)\in\mathrm{AvgBPP}$, and so, by Lemma 6, Item 3, there is a randomized polytime algorithm $B$ such that $B(x,1^{2^{2|x|}})\neq L(x)$ with probability less than $3\cdot2^{-2n}<1/(3\cdot2^n)$, where the probability is over a uniformly random $x\in\{0,1\}^n$ and the internal randomness of $B$. It follows that, for every $x\in\{0,1\}^n$, $B(x,1^{2^{2|x|}})\neq L(x)$ with probability less than $1/3$ over its internal randomness. Hence, the randomized algorithm $B'(x):=B(x,1^{2^{2|x|}})$ decides if $x\in L$ with probability at least $2/3$, for all sufficiently large inputs $x\in\{0,1\}^n$, yielding $L\in\mathrm{BPE}$.\hfill$\triangleleft$
\paragraph{Corollary 12.} If $\mathrm{DistP}^{\Sigma_2^{\mathrm P}}\subseteq\mathrm{AvgBPP}$, then there is an $\varepsilon$-error pseudodeterministic PRG of seed length $O(\log(n/\varepsilon))$.
\paragraph{Proof.} Immediate by combining Lemmas 9–11.\hfill$\triangleleft$
\paragraph{Lemma 23 ($\mathrm{rK}^t$ Reconstruction Lemma).} Assume the existence of a pseudodeterministic PRG. For $\varepsilon>0$, $x\in\{0,1\}^n$, $s\in\mathbb N$, and $k\in\mathbb N$ satisfying $k\leq2n$, let $D$ be a randomized algorithm that takes an advice $\beta$ and runs in time $t_D$ such that $D$ $\varepsilon$-distinguishes $\mathrm{DP}_k(x;U_{nk})$ from $U_{nk+k}$. Then there is a polynomial $p'_{\mathrm{DP}}$ such that
\[\mathrm{rK}^{p'_{\mathrm{DP}}(nt_D/\varepsilon)}(x\mid\beta)\leq k+\log p'_{\mathrm{DP}}(nt_D/\varepsilon).\]
\paragraph{Proof.} By averaging, arguing as in Lemma 22 up to Equation (4),
\[\Pr_{r_D\sim\{0,1\}^s}[\mathbb E_y[D(\mathrm{DP}_k(x,y),r_D)]-\mathbb E_z[D(z,r_D)]\geq\varepsilon/2]\geq\varepsilon/2.\tag{5}\]
That is, with probability at least $\varepsilon/2$ over its internal randomness $r_D$, $D(-,r_D)$ is an $(\varepsilon/2)$-distinguisher for $\mathrm{DP}_k(x;U_{nk})$ and $U_{nk+k}$. Let us say that $r_D$ is good if this is true.
Let $\mathrm{Recon}^{(-)}:\{0,1\}^k\times\{0,1\}^r\to\{0,1\}^n$ be the reconstruction procedure from Lemma 14 that works for distinguishing parameter $\varepsilon/2$. By definition of Recon, if $r_D$ is good, then
\[\Pr_{w\sim\{0,1\}^r}[\mathrm{Recon}^{D(-,r_D)}(A(x,w),w)=x]\geq1/\operatorname{poly}(n/\varepsilon),\]
and so
\[\Pr_{w,r_D}[\mathrm{Recon}^{D(-,r_D)}(A(x,w),w)=x]\geq\frac\varepsilon2\frac1{\operatorname{poly}(n/\varepsilon)}\stackrel{\mathrm{def}}=\varepsilon'.\]
Observe that by Lemma 14, the condition $\mathrm{Recon}^{D(-,r_D)}(A(x,w),w)=x$ can be checked by a circuit of size $s=\operatorname{poly}(nt_D/\varepsilon)$ given $(r_D,w)$ as input. Consider an assumed pseudodeterministic PRG
\[G:\{0,1\}^\ell\to\{0,1\}^{|r_D|+|w|}\]
that $(\varepsilon'/2)$-fools all circuits of size $s$, where $\ell\in O(\log(s/\varepsilon'))\subseteq O(\log(ns/\varepsilon))$. Let $M_G$ be an algorithm as described in Definition 8, running in time $\operatorname{poly}(s/\varepsilon')\leq\operatorname{poly}(nt_D/\varepsilon)$, such that for all $\sigma\in\{0,1\}^\ell$,
\[\Pr_{r_M\sim\{0,1\}^{\operatorname{poly}(nt_D/\varepsilon)}}[M_G(\sigma,r_M)=G(\sigma)]\geq1-2^{-s},\]
where $r_M$ denotes the internal randomness of $M_G$. By definition of $G$,
\[\Pr_\sigma[\mathrm{Recon}^{D(-,G_0(\sigma))}(A(x,G_1(\sigma)),G_1(\sigma))=x]\geq\varepsilon'/2,\]
where $G_0(\sigma)=G(\sigma)_{1\ldots|r_D|}$ (the $|r_D|$-length prefix) and $G_1(\sigma)=G(\sigma)_{|r_D|+1\ldots|r_D|+|w|}$ (the remaining suffix). So there must exist a seed $\sigma\in\{0,1\}^\ell$ such that
\[\mathrm{Recon}^{D(-,G_0(\sigma))}(A(x,G_1(\sigma)),G_1(\sigma))=x.\]
Let $\sigma$ be a seed with this property, and let $\alpha=A(x,G_1(\sigma))\in\{0,1\}^k$. The definition of $M_G$ implies that $\alpha$ and $\sigma$ are such that
\[\Pr_{r_M}[\mathrm{Recon}^{D(-,M_G^0(\sigma,r_M))}(\alpha,M_G^1(\sigma,r_M))=x]\geq1-2^{-s}>2/3,\]
where $M_G^0(\sigma,r_M)$ denotes the $|r_D|$-length prefix of $M_G(\sigma,r_M)$, and $M_G^1(\sigma,r_M)$ the remaining suffix of $M_G(\sigma,r_M)$. It is easy to verify that the reconstruction procedure $\mathrm{Recon}^{D(-,M_G^0(\sigma,r_M))}(\alpha,M_G(\sigma,-))$ runs in time $\operatorname{poly}(nt_D/\varepsilon)$ overall. By definition of $\mathrm{rK}^t$, for some polynomial $p'_{\mathrm{DP}}$, we get
\begin{align*}
\mathrm{rK}^{p'_{\mathrm{DP}}(nt_D/\varepsilon)}(x\mid\beta)&\leq|\alpha|+|\sigma|+O(\log(ns/\varepsilon))\\
&\leq k+\log p'_{\mathrm{DP}}(nt_D/\varepsilon),
\end{align*}
as required.\hfill$\triangleleft$
\paragraph{Lemma 26, original Item 2.} If $\mathrm{DistP}^{\Sigma_2^{\mathrm P}}\subseteq\mathrm{AvgBPP}$, then there exist polynomials $p'$ and $p'_0$ such that for all sufficiently large $n,m\in\mathbb N$, all $t\geq p'_0(n,m)$, and all $0\leq\tau\leq t$ the following holds: for every $x\in\{0,1\}^n$ and every family $\{y_r\in\{0,1\}^m\mid r\in\{0,1\}^\tau\}$, with probability at least $9/10$ over $r\in\{0,1\}^\tau$,
\[\mathrm{rK}^t(x,y_r\mid r)>\mathrm{rK}^{p'(t)}(x\mid r)+\mathrm{rK}^{p'(t)}(y_r\mid x,r)-\log p'(t).\]
Moreover, the special case without conditioning on $r$ also holds for both items above, i.e., when $r$ is the empty string (for $\tau=0$) and $y$ is an arbitrary single string.
\paragraph{Proof of Item 2.} Argue as in the proof of Item 1, but at Equation (9), apply Lemma 23 instead of Lemma 22 to get
\[\mathrm{rK}^{q(t)}(x\mid r)\leq k+\log q(t)\]
for some polynomial $q$ such that $q(t)\geq p'_{\mathrm{DP}}(t_B\cdot3\cdot n)$, where $p'_{\mathrm{DP}}$ is the polynomial from Lemma 23 and $t_B$ denotes the time required to compute $B(-,v,w,r,1^s)$. We then define $k:=\mathrm{rK}^{q(t)}(x\mid r)-\log q(t)-1$ so the above equation does not hold. Similarly, at Equation (11), apply Lemma 23 to get
\[\mathrm{rK}^{q'(t)}(y_r\mid x,r)\leq k'+\log q'(t)\]
for some polynomial $q'$ with $q'(t)\geq p'_{\mathrm{DP}}(t_{B'}\cdot8\cdot m)$, where $t_{B'}$ denotes the time required to compute the distinguisher $B'$ defined as above. Then define $k':=\mathrm{rK}^{q'(t)}(y_r\mid x,r)-\log q'(t)-1$.
Following the proof of Item 1 up to Equation (12), we get that
\[\mathrm{pK}^t(x,y_r\mid r)>k+k'-d\log t-\log n.\]
The definitions of $k,k'$, and $\mathrm{rK}^t$ then imply that
\begin{multline*}
\mathrm{rK}^t(x,y_r\mid r)>\mathrm{rK}^{q(t)}(x\mid r)+\mathrm{rK}^{q'(t)}(y_r\mid x,r)\\
-\log q(t)-\log q'(t)-d\log t-12,
\end{multline*}
and so
\[\mathrm{rK}^t(x,y_r\mid r)>\mathrm{rK}^{p(t)}(x\mid r)+\mathrm{rK}^{p(t)}(y_r\mid x,r)-\log p(t)\]
for the polynomial $p(t):=(t^{d+1})\cdot q(t)\cdot q'(t)$.\hfill$\triangleleft$
\paragraph{Lemma 33, original Item 2.} Assume in addition that $\mathrm{DistP}^{\Sigma_2^{\mathrm P}}\subseteq\mathrm{AvgBPP}$ (which holds in particular for $\ell\geq2$). Then there exist polynomials $q'$ and $q'_0$ such that, for all sufficiently large $n\in\mathbb N$, all $x\in L_n$, and all $t\geq q'_0(n)$, we have with probability at least $\delta-1/10$ over $r\in\{0,1\}^{n^c}$ that an $L$-witness $y_{x,r}$ exists, and
\[\mathrm{rK}^{q'(t)}(x,y_{x,r}\mid r)\leq\mathrm{rK}^t(x\mid r)+\log q'(t).\]
Moreover, for the case of no BP quantifier, i.e., for $L\in\Sigma_\ell^{\mathrm P}$, we get the same conclusions (with probability $1$) without conditioning on $r$.
\paragraph{Proof of Item 2.} The proof is similar to that of Item 1, except at Equation (19), we observe that for all $t\geq q_0(n)$,
\[\mathrm{pK}^{2t,\Sigma_\ell^{\mathrm P}}(x,y_{x,r}\mid r)\leq\mathrm{pK}^t(x\mid r)+d\log n\leq\mathrm{rK}^t(x\mid r)+d\log n,\]
and then define $s(x,r):=\mathrm{rK}^t(x\mid r)+d\log n$. We then follow the proof of Item 1 up to Equation (21), where we apply Lemma 23 instead of Lemma 22 to obtain
\[\mathrm{rK}^{p''(t)}(x,y_{x,r}\mid r,s,n,t)\leq\mathrm{rK}^t(x\mid r)+d\log n+\log p''(t)+10\]
for some polynomial $p''$ with $p''(t)\geq p'_{\mathrm{DP}}(t_B\cdot3\cdot(n+n^c))$, where $p'_{\mathrm{DP}}$ is the polynomial from Lemma 23 and $t_B$ denotes the time required to compute $B(-,U_{2t},r,1^k,1^s,1^n)$. It follows that $\mathrm{rK}^{q'(t)}(x,y_{x,r}\mid r)\leq\mathrm{rK}^t(x\mid r)+\log q'(t)$ for the polynomial $q'(t):=t^{d+4}\cdot p''(t)$.\hfill$\triangleleft$
\paragraph{Lemma 34, original Item 2.} Assume in addition that $\mathrm{DistP}^{\Sigma_2^{\mathrm P}}\subseteq\mathrm{AvgBPP}$ (which holds in particular when $\ell\geq2$). Let $p'_0,p'$ and $q'_0,q'$ be the polynomials from Items 2 of Lemmas 26 and 33 respectively. Then for every $t\geq\max\{p'_0(n),q'_0(n+n^c)\}$, with probability at least $\delta-1/5$ over $r$, an $L$-witness $y_{x,r}$ exists, and
\[\mathrm{rK}^{p'(q'(t))}(y_{x,r}\mid x,r)\leq(\mathrm{rK}^t(x\mid r)-\mathrm{rK}^{p'(q'(t))}(x\mid r))+2\log p'(q'(t)).\]
Moreover, for $L\in\Sigma_\ell^{\mathrm P}$, we get the same conclusions without conditioning on $r$.
The proof of Item 2 is the same, except we apply Items 2 of Lemmas 33 and 26 respectively.
